%2multibyte Version: 5.50.0.2960 CodePage: 65001
% MQF_Manuale_Master_SW55_Memoir.tex
% Master del manuale MQF costruito sullo shell SW5.5 "Memoir Class for Configurable Typesetting".
% Versione modulare: i capitoli scritti sono inclusi con \input.%
% POSIZIONE DEI FILE
% Questo master deve stare nella cartella 01_Manuale.
% I capitoli devono stare nella sottocartella 01_Manuale/Capitoli.
% Non mettere il master dentro Capitoli.
% I capitoli futuri restano come traccia, ma non generano pagine nel documento.
% Evita pagine bianche dovute all'apertura dei capitoli su pagina dispari.
% ------------------------------------------------------------
% Nomi italiani
% ------------------------------------------------------------
% ------------------------------------------------------------
% Macro matematiche ufficiali del progetto MQF
% Coerenti con MQF_Notazione.tex
% ------------------------------------------------------------
% ------------------------------------------------------------
% Ambienti di tipo teorema.
% I nomi interni restano compatibili con lo shell SW.
% Le etichette stampate sono in italiano.
% ------------------------------------------------------------
% MQF_Manuale_Master_SW55_Memoir.tex
% Master del manuale MQF costruito sullo shell SW5.5 "Memoir Class for Configurable Typesetting".
% Versione modulare: i capitoli scritti sono inclusi con \input.%
% POSIZIONE DEI FILE
% Questo master deve stare nella cartella 01_Manuale.
% I capitoli devono stare nella sottocartella 01_Manuale/Capitoli.
% Non mettere il master dentro Capitoli.
% I capitoli futuri restano come traccia, ma non generano pagine nel documento.
% Evita pagine bianche dovute all'apertura dei capitoli su pagina dispari.
% ------------------------------------------------------------
% Nomi italiani
% ------------------------------------------------------------
% ------------------------------------------------------------
% Macro matematiche ufficiali del progetto MQF
% Coerenti con MQF_Notazione.tex
% ------------------------------------------------------------
% ------------------------------------------------------------
% Ambienti di tipo teorema.
% I nomi interni restano compatibili con lo shell SW.
% Le etichette stampate sono in italiano.
% ------------------------------------------------------------

%TCIDATA{OutputFilter=latex2.dll}
%TCIDATA{Version=5.50.0.2960}
%TCIDATA{LaTeXparent=1,1,MQF_Manuale_Master.tex}


\chapter{Elementi di teoria della probabilit\`{a}}

\label{cap:probabilita}

\section{Obiettivi della lezione}

La finanza quantitativa studia grandezze economiche e finanziarie che non sono
note con certezza nel momento in cui una decisione viene assunta. Il
rendimento futuro di un titolo, il valore di un portafoglio alla fine di un
periodo, l'eventuale default di un debitore, la variazione di un tasso di
interesse o il superamento di una soglia di perdita sono esempi di fenomeni il
cui esito dipende da condizioni future non pienamente osservabili al tempo
della decisione.

Il primo obiettivo di questo capitolo \`{e} introdurre il linguaggio
matematico di base con cui tale incertezza viene rappresentata. Questo
linguaggio \`{e} quello della teoria della probabilit\`{a}. In particolare, il
capitolo introduce gli elementi costitutivi di uno spazio di probabilit\`{a}:
\[
(\Omega,\mathcal{F},\mathbb{P}).
\]
L'insieme $\Omega$ rappresenta gli esiti elementari possibili, $\mathcal{F}$
rappresenta la famiglia degli eventi ai quali il modello assegna significato
probabilistico, mentre $\mathbb{P}$ assegna a ciascun evento una probabilit\`{a}.

Il secondo obiettivo \`{e} chiarire come si calcolano le probabilit\`{a} di
eventi composti e come si interpreta il condizionamento. In finanza, il
condizionamento \`{e} essenziale: un segnale di mercato, un deterioramento del
rating, una variazione dello spread o un nuovo dato macroeconomico modificano
l'informazione disponibile e quindi la valutazione probabilistica degli eventi futuri.

Al termine del capitolo lo studente dovr\`{a} essere in grado di:

\begin{itemize}
\item distinguere tra esiti elementari, eventi e famiglie di eventi;

\item interpretare uno spazio di probabilit\`{a} in un contesto finanziario;

\item utilizzare le operazioni fondamentali sugli eventi;

\item applicare le regole elementari della probabilit\`{a};

\item calcolare probabilit\`{a} condizionate;

\item distinguere tra eventi disgiunti ed eventi indipendenti;

\item applicare la formula delle probabilit\`{a} totali e la formula di Bayes;

\item interpretare l'aggiornamento bayesiano come revisione della
probabilit\`{a} in presenza di nuova informazione.
\end{itemize}

Il capitolo prepara il passaggio al Capitolo 2. In questa fase l'attenzione
\`{e} posta sugli eventi. Nel capitolo successivo gli scenari saranno
trasformati in grandezze numeriche mediante il concetto di variabile casuale.

\section{Motivazione finanziaria: incertezza, scenari ed eventi}

Consideriamo un investitore che osserva oggi il prezzo di un titolo e deve
valutare il possibile valore del proprio investimento alla fine del periodo
successivo. Al tempo iniziale, indicato con $t=0$, il valore futuro non \`{e}
noto. Esistono diversi scenari compatibili con l'informazione disponibile: il
prezzo potrebbe aumentare, rimanere sostanzialmente stabile oppure diminuire.

In termini elementari, possiamo rappresentare questi scenari mediante uno
spazio campionario finito:
\[
\Omega=\{\omega_{1},\omega_{2},\omega_{3}\},
\]
dove $\omega_{1}$ rappresenta uno scenario di ribasso, $\omega_{2}$ uno
scenario di stabilit\`{a} e $\omega_{3}$ uno scenario di rialzo. Lo spazio
$\Omega$ non \`{e} ancora un modello probabilistico completo: esso elenca
soltanto gli esiti possibili.

Se definiamo
\[
A=\{\omega_{1}\},
\]
l'evento $A$ pu\`{o} essere interpretato come evento di perdita. Se definiamo
\[
B=\{\omega_{2},\omega_{3}\},
\]
l'evento $B$ rappresenta la circostanza in cui il rendimento non \`{e}
negativo. L'evento certo \`{e} $\Omega$, mentre l'evento impossibile \`{e}
$\emptyset$.

Un secondo esempio riguarda il rischio di credito. Un debitore pu\`{o}
rimborsare regolarmente oppure andare in default. In una rappresentazione
elementare si pu\`{o} porre
\[
\Omega=\{\omega_{N},\omega_{D}\},
\]
dove $\omega_{N}$ indica lo stato di non default e $\omega_{D}$ indica lo
stato di default. L'evento di default \`{e}
\[
D=\{\omega_{D}\}.
\]
La probabilit\`{a} $\mathbb{P}(D)$ rappresenta la probabilit\`{a} di default
sull'orizzonte temporale considerato. Tale probabilit\`{a}, tuttavia, pu\`{o}
essere rivista se si osserva un nuovo segnale informativo, ad esempio un
peggioramento del rating o un aumento dello spread.

Questi esempi mostrano che un modello probabilistico richiede tre elementi:
\[
\text{scenari possibili}\quad\longrightarrow\quad\text{eventi rilevanti}%
\quad\longrightarrow\quad\text{probabilit\`{a} degli eventi}.
\]
La tripla $(\Omega,\mathcal{F},\mathbb{P})$ formalizza precisamente questi tre elementi.

\section{Spazio campionario ed eventi}

\begin{definition}
Lo spazio campionario $\Omega$ \`{e} l'insieme degli esiti elementari
possibili di un esperimento aleatorio. Un elemento $\omega\in\Omega$ \`{e}
detto evento o esito elementare.
\end{definition}

Nel contesto finanziario, gli esiti elementari possono rappresentare scenari
di mercato, stati di credito, evoluzioni di un tasso di interesse oppure
possibili realizzazioni di un insieme di fattori di rischio.

\begin{definition}
Un evento \`{e} un sottoinsieme dello spazio campionario. Se $A\subseteq
\Omega$, dire che l'evento $A$ si verifica significa che l'esito elementare
realizzato appartiene ad $A$.
\end{definition}

Nel caso di uno spazio finito, ad esempio
\[
\Omega=\{\omega_{1},\omega_{2},\omega_{3},\omega_{4}\},
\]
un evento pu\`{o} essere
\[
A=\{\omega_{1},\omega_{2}\}.
\]
Se $\omega_{1}$ o $\omega_{2}$ si realizza, allora si dice che $A$ si \`{e}
verificato. Se invece si realizza $\omega_{3}$ o $\omega_{4}$, l'evento $A$
non si \`{e} verificato.

In questo capitolo useremo soprattutto spazi campionari finiti o discreti. Gli
spazi continui saranno richiamati solo quando necessario e saranno ripresi nel
Capitolo 2, quando introdurremo variabili casuali e distribuzioni.

\section{Spazi campionari generati da esperimenti ripetuti}

Molti problemi finanziari non riguardano un singolo esperimento aleatorio, ma
una sequenza di esperimenti osservati nel tempo. Si pensi, ad esempio, ai
movimenti giornalieri di un prezzo, all'evoluzione periodica del rating di una
controparte, alla successione di rialzi e ribassi del valore di un
portafoglio, oppure alla simulazione di una traiettoria di mercato su pi\`{u}
date future.

Supponiamo che l'incertezza sia osservata in $n$ date successive e che, alla
data $i=1,\ldots,n$, gli esiti possibili siano descritti da uno spazio
campionario $\Omega_{i}$. Lo spazio campionario complessivo \`{e} allora il
prodotto cartesiano
\[
\Omega=\Omega_{1}\times\Omega_{2}\times\cdots\times\Omega_{n}.
\]
Un esito elementare $\omega\in\Omega$ \`{e} una $n$-pla
\[
\omega=(\omega_{1},\omega_{2},\ldots,\omega_{n}),
\]
dove $\omega_{i}\in\Omega_{i}$ rappresenta l'esito elementare realizzato nella
$i$-esima osservazione.

Questa notazione \`{e} particolarmente importante per il seguito del manuale.
Essa permette di interpretare un esito elementare non come un singolo stato
isolato ad una sola data $i$, ma come una \textit{sequenza ordinata} di
realizzazioni. In questo senso, un elemento $\omega\in\Omega$ pu\`{o} essere
letto come una \textit{traiettoria} elementare.

Nel caso in cui lo stesso esperimento sia ripetuto $n$ volte e ciascuno spazio
degli esiti sia lo stesso a ogni data, cio\`{e}
\[
\Omega_{1}=\Omega_{2}=\cdots=\Omega_{n}=\Omega^{\ast},
\]
si scrive pi\`{u} semplicemente
\[
\Omega=\underset{\text{n volte}}{\underbrace{\Omega^{\ast}\times\Omega^{\ast
}\times...\times\Omega^{\ast}}}=(\Omega^{\ast})^{n}.
\]
Se $\Omega^{\ast}$ contiene $m$ esiti elementari, allora lo spazio prodotto
$(\Omega^{\ast})^{n}$ contiene
\[
m^{n}%
\]
esiti elementari.

\begin{example}
Consideriamo tre movimenti successivi del prezzo di un titolo. A ogni data il
prezzo pu\`{o} registrare un rialzo, indicato con $u$, oppure un ribasso,
indicato con $d$. Lo spazio degli esiti di una singola osservazione \`{e}
\[
\Omega^{\ast}=\{u,d\}.
\]
Per tre osservazioni successive, lo spazio campionario \`{e}
\[
\Omega=(\Omega^{\ast})^{3}=\{u,d\}^{3}.
\]
Gli esiti elementari sono le otto sequenze
\[
\Omega=\{uuu,uud,udu,udd,duu,dud,ddu,ddd\}.
\]
Ad esempio,
\[
\omega=udd
\]
rappresenta la traiettoria in cui il primo movimento \`{e} un rialzo, mentre
il secondo e il terzo sono ribassi.
\end{example}

In questa rappresentazione, gli eventi sono sottoinsiemi di sequenze. Ad
esempio, nell'esempio precedente, l'evento
\[
A=\{uuu,uud,udu,udd\}
\]
rappresenta la circostanza in cui il primo movimento \`{e} un rialzo.
L'evento
\[
B=\{udd,dud,ddu\}
\]
rappresenta invece la circostanza in cui si osserva esattamente un rialzo nei
tre periodi. L'evento
\[
C=\{uuu,uud,udu,duu\}
\]
rappresenta la circostanza in cui si osservano almeno due rialzi.

Questi esempi mostrano che, quando gli esiti elementari sono sequenze, gli
eventi possono descrivere propriet\`{a} della traiettoria: il primo movimento,
il numero totale di rialzi, il verificarsi di due ribassi consecutivi, oppure
il superamento di una soglia lungo il percorso.

\begin{example}
Consideriamo ora quattro date trimestrali di osservazione del merito
creditizio di una controparte. A ogni trimestre si osserva se il merito
creditizio rimane stabile, indicato con $s$, oppure si deteriora, indicato con
$d$. Lo spazio degli esiti di una singola osservazione \`{e}
\[
\Omega^{\ast}=\{s,d\}.
\]
Su quattro trimestri,
\[
\Omega=(\Omega^{\ast})^{4}=\{s,d\}^{4},
\]
e quindi lo spazio campionario contiene
\[
2^{4}=16
\]
traiettorie elementari.

L'esito
\[
\omega=sdds
\]
rappresenta una traiettoria nella quale il merito creditizio resta stabile nel
primo trimestre, si deteriora nel secondo e nel terzo, e torna stabile nel quarto.

L'evento
\[
A=\{\omega\in\Omega:\text{ in }\omega\text{ compaiono almeno due simboli }d\}
\]
rappresenta la circostanza in cui si osservano almeno due deterioramenti nel
periodo considerato. Questo evento non riguarda un singolo trimestre, ma una
propriet\`{a} dell'intera traiettoria.
\end{example}

La stessa logica si estende al caso di un numero infinito di osservazioni. Se
l'esperimento viene ripetuto indefinitamente e lo spazio elementare \`{e}
$\Omega^{\ast}$, si pu\`{o} considerare lo spazio
\[
\Omega=(\Omega^{\ast})^{\mathbb{N}}.
\]
In questo caso un esito elementare \`{e} una successione infinita
\[
\omega=(\omega_{1},\omega_{2},\ldots).
\]
Nel presente manuale lavoreremo prevalentemente con orizzonti finiti
$t=0,1,\ldots,T$, ma la notazione per successioni infinite \`{e} utile per
comprendere il legame con modelli stocastici pi\`{u} generali.

Questa costruzione sar\`{a} ripresa pi\`{u} volte. Nei processi stocastici,
una traiettoria descrive l'evoluzione temporale di una variabile aleatoria.
Negli alberi binomiali, ogni percorso dalla radice, detta \textit{foglia},
rappresenta una sequenza di rialzi e ribassi. Nella simulazione, ogni replica
genera una traiettoria casuale. Nella programmazione stocastica, uno scenario
pu\`{o} essere interpretato come una realizzazione coerente dell'incertezza
lungo l'orizzonte temporale considerato.

\section{Operazioni sugli eventi}

Siano $A,B\subseteq\Omega$ due eventi.

L'\textit{unione}
\[
A\cup B
\]
\`{e} l'evento che si verifica quando si verifica almeno uno tra $A$ e $B$.
L'\textit{intersezione}
\[
A\cap B
\]
\`{e} l'evento che si verifica quando si verificano simultaneamente $A$ e $B$.
Il complemento
\[
A^{c}=\Omega\setminus A
\]
\`{e} l'evento che si verifica quando $A$ non si verifica.

Due eventi $A$ e $B$ si dicono disgiunti se
\[
A\cap B=\emptyset.
\]
In questo caso non possono verificarsi simultaneamente.

\begin{example}
Supponiamo che $\Omega=\{\omega_{1},\omega_{2},\omega_{3}\}$ rappresenti tre
scenari di mercato: ribasso, stabilit\`{a} e rialzo. Se
\[
A=\{\omega_{1}\},\qquad B=\{\omega_{1},\omega_{2}\},
\]
allora $A$ rappresenta l'evento di perdita, mentre $B$ rappresenta l'evento in
cui non vi \`{e} rialzo. Si ha
\[
A\cap B=\{\omega_{1}\},\qquad A\cup B=\{\omega_{1},\omega_{2}\},\qquad
A^{c}=\{\omega_{2},\omega_{3}\}.
\]

\end{example}

\subsection{Famiglie di eventi}
In molte situazioni non interessa considerare un singolo evento, ma un insieme di eventi accomunati da una certa propriet\`{a}. 
Una famiglia di eventi \`{e} un insieme del tipo
\[
\mathcal{A}=\{A_i:i\in I\},
\]
dove ogni \(A_i\) \`{e} un evento.
Occorre distinguere accuratamente la famiglia \(\mathcal{A}\) dall'evento ottenuto unendo gli eventi che la compongono:
\[
\bigcup_{i\in I} A_i.
\]
La famiglia \(\mathcal{A}\) \`{e} un insieme i cui elementi sono eventi, ma non è un evento; l'unione \(\bigcup_{i\in I}A_i\), invece, \`{e} un sottoinsieme di \(\Omega\), e quindi pu\`{o} essere un evento.

\begin{example}
Sia
\[
\Omega=\{1,2,3,4,5,6\}
\]
lo spazio dei risultati del lancio di un dado. Consideriamo gli eventi
\[
A=\{2,4,6\},\qquad B=\{1,3,5\}.
\]
La famiglia
\[
\mathcal{A}=\{A,B\}
\]
ha come elementi due eventi: l'evento ``esce un numero pari'' e l'evento ``esce un numero dispari''. 
Essa non coincide con l'evento
\[
A\cup B=\{1,2,3,4,5,6\}=\Omega.
\]
Infatti \(\mathcal{A}\) \`{e} un insieme di sottoinsiemi di \(\Omega\), mentre \(A\cup B\) \`{e} un sottoinsieme di \(\Omega\). 
La probabilit\`{a} \(P(A\cup B)\) \`{e} ben definita; al contrario, \(P(\mathcal{A})\) non ha significato nello spazio di probabilit\`{a} originario, perch\'{e} \(\mathcal{A}\) non \`{e} un evento di \(\mathcal{F}\), ma una collezione di eventi.
\end{example}

In altre parole:
\[
\mathcal{A}=\{A,B\}\neq A\cup B.
\]
e
\[
\mathcal{A}=\{A_i:i\in I\},
\qquad
A_i\subseteq\Omega,
\qquad
\bigcup_{i\in I}A_i\subseteq\Omega.
\]
Le famiglie di eventi permettono di descrivere in modo compatto classi di eventi, operazioni collettive e strutture logiche. 
Per esempio, data una famiglia \(\{A_i:i\in I\}\), l'evento
\[
\bigcup_{i\in I}A_i
\]
rappresenta il verificarsi di almeno uno degli eventi \(A_i\), mentre
\[
\bigcap_{i\in I}A_i
\]
rappresenta il verificarsi simultaneo di tutti gli eventi della famiglia. 
Analogamente,
\[
\left(\bigcup_{i\in I}A_i\right)^c
\]
rappresenta l'evento in cui nessuno degli eventi \(A_i\) si verifica.

\subsection{Partizione di $\Omega$}
Una importante famiglia di eventi è la \textbf{partizione di $\Omega$}.

\begin{definition}
Una famiglia di eventi $B_{1},\ldots,B_{n}$ \`{e} una partizione di $\Omega$
se gli eventi sono due a due disgiunti e la loro unione \`{e} l'intero spazio
campionario:
\[
B_{i}\cap B_{j}=\emptyset\quad(i\neq j),\qquad\bigcup_{i=1}^{n}B_{i}=\Omega.
\]

\end{definition}

Le partizioni sono importanti perch\'{e} permettono di decomporre un problema
probabilistico rispetto a scenari alternativi e mutuamente esclusivi.

\section{Sigma-algebra e informazione}

\begin{definition}
Una sigma-algebra $\mathcal{F}$ su $\Omega$ \`{e} una famiglia di sottoinsiemi
di $\Omega$ tale che:

\begin{enumerate}
\item $\Omega\in\mathcal{F}$;

\item se $A\in\mathcal{F}$, allora $A^{c}\in\mathcal{F}$;

\item se $A_{1},A_{2},\ldots$ appartengono a $\mathcal{F}$, allora
$\bigcup_{n=1}^{\infty} A_{n}\in\mathcal{F}$.
\end{enumerate}
\end{definition}

Nel seguito, gli elementi di $\mathcal{F}$ saranno chiamati eventi. In uno
spazio campionario finito si lavora spesso con l'insieme delle parti di
$\Omega$; in tal caso tutti i sottoinsiemi di $\Omega$ sono eventi. Per il
manuale \`{e} tuttavia utile mantenere la notazione generale $\mathcal{F}$,
perch\'{e} essa consente di rappresentare formalmente l'informazione disponibile.

L'interpretazione finanziaria \`{e} la seguente: $\mathcal{F}$ descrive quali
eventi sono distinguibili nel modello. Una sigma-algebra pi\`{u} ricca
consente di distinguere pi\`{u} eventi; una sigma-algebra pi\`{u} povera
rappresenta un livello informativo pi\`{u} aggregato.

\begin{example}
Consideriamo tre movimenti successivi di un titolo, ciascuno dei quali pu\`{o}
essere un rialzo $u$ oppure un ribasso $d$. Lo spazio campionario \`{e}
\[
\Omega=\{uuu,uud,udu,udd,duu,dud,ddu,ddd\}.
\]
Prima di osservare i movimenti, l'informazione \`{e} minima. Dopo il primo
movimento si pu\`{o} distinguere tra gli scenari che iniziano con $u$ e quelli
che iniziano con $d$. Dopo due movimenti, la distinzione diventa pi\`{u} fine.
Dopo tre movimenti l'informazione \`{e} completa. Questa successione di
livelli informativi anticipa il ruolo delle filtrazioni $\{\mathcal{F}%
_{t}\}_{t=0}^{T}$, che saranno introdotte nei capitoli successivi.
\end{example}

\section{Misura di probabilit\`{a}}

\begin{definition}
Una misura di probabilit\`{a} su $(\Omega,\mathcal{F})$ \`{e} una funzione
\[
\mathbb{P}:\mathcal{F}\longrightarrow\lbrack0,1]
\]
tale che:

\begin{enumerate}
\item $\mathbb{P}(A)\geq0$ per ogni $A\in\mathcal{F}$;

\item $\mathbb{P}(\Omega)=1$;

\item se $A_{1},A_{2},\ldots$ sono eventi due a due disgiunti, allora
\[
\mathbb{P}\left(  \bigcup_{n=1}^{\infty} A_{n}\right)  = \sum_{n=1}^{\infty
}\mathbb{P}(A_{n}).
\]

\end{enumerate}
\end{definition}

Da questi assiomi discendono alcune regole di calcolo elementari. Per ogni
$A,B\in\mathcal{F}$,
\[
\mathbb{P}(A^{c})=1-\mathbb{P}(A),
\]
e
\[
\mathbb{P}(A\cup B)=\mathbb{P}(A)+\mathbb{P}(B)-\mathbb{P}(A\cap B).
\]
Se $A$ e $B$ sono disgiunti, allora
\[
\mathbb{P}(A\cup B)=\mathbb{P}(A)+\mathbb{P}(B).
\]


\begin{example}
Supponiamo che uno spazio campionario finito sia dato da
\[
\Omega=\{\omega_{1},\omega_{2},\omega_{3}\},
\]
con
\[
\mathbb{P}(\{\omega_{1}\})=0.25,\qquad\mathbb{P}(\{\omega_{2}\})=0.35,\qquad
\mathbb{P}(\{\omega_{3}\})=0.40.
\]
Se $A=\{\omega_{1}\}$ e $B=\{\omega_{2},\omega_{3}\}$, allora
\[
\mathbb{P}(A)=0.25,\qquad\mathbb{P}(B)=0.35+0.40=0.75.
\]

\end{example}

\section{Spazio di probabilit\`{a}}

\begin{definition}
Uno spazio di probabilit\`{a} \`{e} una tripla
\[
(\Omega,\mathcal{F},\mathbb{P}),
\]
dove $\Omega$ \`{e} lo spazio campionario, $\mathcal{F}$ \`{e} una
sigma-algebra di eventi e $\mathbb{P}$ \`{e} una misura di probabilit\`{a} su
$\mathcal{F}$.
\end{definition}

Questa tripla separa tre piani concettuali distinti: gli scenari possibili,
gli eventi osservabili o modellizzati, e la probabilit\`{a} assegnata a tali
eventi. La distinzione \`{e} essenziale in finanza, perch\'{e} una stessa
struttura di scenari pu\`{o} essere valutata con misure di probabilit\`{a}
diverse, ad esempio una misura storica oppure, nei modelli di pricing, una
misura risk-neutral $\mathbb{Q}$.

\section{Probabilit\`{a} condizionata}

\begin{definition}
Siano $A,B\in\mathcal{F}$ e sia $\mathbb{P}(B)>0$. La probabilit\`{a}
condizionata di $A$ dato $B$ \`{e} definita da
\[
\mathbb{P}(A\mid B)=\frac{\mathbb{P}(A\cap B)}{\mathbb{P}(B)}.
\]

\end{definition}

La probabilit\`{a} condizionata rappresenta la probabilit\`{a} dell'evento $A$
quando si assume noto che l'evento $B$ si sia verificato. In questa
prospettiva, $B$ restringe lo spazio informativo rilevante.

Dalla definizione segue la regola del prodotto:
\[
\mathbb{P}(A\cap B)=\mathbb{P}(A\mid B)\mathbb{P}(B).
\]
Simmetricamente, se $\mathbb{P}(A)>0$,
\[
\mathbb{P}(A\cap B)=\mathbb{P}(B\mid A)\mathbb{P}(A).
\]


\begin{example}
Questo esempio ha funzione qualitativa. Sia $D$ l'evento di default di un
debitore e sia $S$ l'evento che rappresenta l'osservazione di un segnale
negativo di mercato. La probabilit\`{a}
\[
\mathbb{P}(D\mid S)
\]
misura la probabilit\`{a} di default dopo avere osservato il segnale negativo.

In generale, $\mathbb{P}(D\mid S)$ pu\`{o} essere molto diversa dalla
probabilit\`{a} non condizionata $\mathbb{P}(D)$. Il segnale $S$ modifica
infatti l'informazione disponibile e restringe l'attenzione agli scenari nei
quali tale segnale si \`{e} verificato.
\end{example}

\begin{example}
Consideriamo due esposizioni creditizie appartenenti allo stesso portafoglio.
Per ciascuna esposizione si osserva, su un dato orizzonte temporale, se vi
\`{e} stato un deterioramento del merito creditizio oppure no. Indichiamo con
$D$ il deterioramento e con $N$ l'assenza di deterioramento. Supponiamo, in
forma puramente stilizzata, che i quattro esiti elementari
\[
\Omega=\{NN,ND,DN,DD\}
\]
siano equiprobabili.

Sia
\[
A=\{DD\}
\]
l'evento in cui entrambe le esposizioni si deteriorano, e sia
\[
B=\{ND,DN,DD\}
\]
l'evento in cui almeno una delle due esposizioni si deteriora. Vogliamo
calcolare la probabilit\`{a} che entrambe le esposizioni si deteriorino
sapendo che almeno una si \`{e} deteriorata.

Per definizione,
\[
\mathbb{P}(A\mid B)=\frac{\mathbb{P}(A\cap B)}{\mathbb{P}(B)}.
\]
Poich\'{e} $A\subseteq B$, si ha $A\cap B=A$. Inoltre,
\[
\mathbb{P}(A)=\frac{1}{4},\qquad\mathbb{P}(B)=\frac{3}{4}.
\]
Pertanto
\[
\mathbb{P}(A\mid B)=\frac{1/4}{3/4}=\frac{1}{3}.
\]


Il risultato mostra che l'informazione \textquotedblleft almeno una
esposizione si \`{e} deteriorata\textquotedblright\ modifica lo spazio degli
scenari rilevanti. Non si considera pi\`{u} l'intero insieme $\Omega$, ma solo
il sottoinsieme $B$ degli scenari compatibili con l'informazione osservata.
\end{example}

\begin{example}
Un modello preliminare assegna a una controparte uno score intero compreso tra
$1$ e $10$, dove valori pi\`{u} elevati indicano migliore qualit\`{a}
creditizia. In assenza di ulteriori informazioni, supponiamo che i dieci
valori siano equiprobabili:
\[
\Omega=\{1,2,\ldots,10\},\qquad\mathbb{P}(\{\omega\})=\frac{1}{10}%
\quad\text{per ogni }\omega\in\Omega.
\]


Sia
\[
A=\{10\}
\]
l'evento in cui la controparte riceve lo score massimo. Supponiamo ora di
sapere che la controparte ha superato una prima soglia di ammissibilit\`{a},
cio\`{e} che lo score \`{e} almeno pari a $5$. L'evento informativo \`{e}
\[
B=\{5,6,7,8,9,10\}.
\]
Vogliamo calcolare
\[
\mathbb{P}(A\mid B).
\]


Poich\'{e} $A\subseteq B$, si ha $A\cap B=A$. Inoltre,
\[
\mathbb{P}(A)=\frac{1}{10},\qquad\mathbb{P}(B)=\frac{6}{10}.
\]
Quindi
\[
\mathbb{P}(A\mid B)=\frac{\mathbb{P}(A\cap B)}{\mathbb{P}(B)}=\frac
{1/10}{6/10}=\frac{1}{6}.
\]


L'informazione che lo score \`{e} almeno pari a $5$ elimina gli scenari
$\{1,2,3,4\}$ dallo spazio rilevante. La probabilit\`{a} condizionata dello
score massimo non \`{e} quindi $1/10$, ma $1/6$, perch\'{e} il calcolo viene
effettuato all'interno del sottoinsieme degli scenari compatibili con
l'informazione disponibile.
\end{example}

I due esempi precedenti mostrano che il condizionamento non modifica gli
eventi in s\`{e}, ma modifica lo spazio informativo rispetto al quale le loro
probabilit\`{a} vengono valutate. In entrambi i casi, l'evento condizionante
seleziona un sottoinsieme di scenari compatibili con l'informazione osservata;
la probabilit\`{a} condizionata \`{e} quindi una probabilit\`{a} relativa a
tale sottoinsieme.

\section{Indipendenza tra eventi}

\begin{definition}
Due eventi $A,B\in\mathcal{F}$ si dicono indipendenti se
\[
\mathbb{P}(A\cap B)=\mathbb{P}(A)\mathbb{P}(B).
\]

\end{definition}

Se $\mathbb{P}(B)>0$, l'indipendenza \`{e} equivalente a
\[
\mathbb{P}(A\mid B)=\mathbb{P}(A).
\]
Quindi $A$ e $B$ sono indipendenti quando la conoscenza del verificarsi di $B$
non modifica la probabilit\`{a} attribuita ad $A$.

\`{e} importante distinguere indipendenza e incompatibilit\`{a}. Se $A$ e $B$
sono disgiunti e hanno probabilit\`{a} positiva, allora non sono indipendenti:
infatti
\[
\mathbb{P}(A\cap B)=0,
\]
mentre
\[
\mathbb{P}(A)\mathbb{P}(B)>0.
\]
In finanza, l'ipotesi di indipendenza deve essere utilizzata con cautela.
Eventi di default apparentemente separati possono diventare fortemente
dipendenti in presenza di uno shock macroeconomico comune.

\section{Formula delle probabilit\`{a} totali}

Sia $B_{1},\ldots,B_{n}$ una partizione di $\Omega$ con $\mathbb{P}(B_{i})>0$
per ogni $i$. Per ogni evento $A\in\mathcal{F}$ vale
\[
\mathbb{P}(A)=\sum_{i=1}^{n}\mathbb{P}(A\mid B_{i})\mathbb{P}(B_{i}).
\]


La formula delle probabilit\`{a} totali consente di calcolare la
probabilit\`{a} di un evento aggregando le probabilit\`{a} condizionate
rispetto a scenari alternativi.

\begin{example}
Supponiamo che l'economia possa trovarsi in tre stati:
\[
B_{1}=\text{crescita},\qquad B_{2}=\text{stagnazione},\qquad B_{3}%
=\text{recessione}.
\]
Sia $D$ l'evento di default di un debitore. Se conosciamo le probabilit\`{a}
$\mathbb{P}(B_{i})$ e le probabilit\`{a} condizionate $\mathbb{P}(D\mid
B_{i})$, allora
\[
\mathbb{P}(D)=\mathbb{P}(D\mid B_{1})\mathbb{P}(B_{1})+\mathbb{P}(D\mid
B_{2})\mathbb{P}(B_{2})+\mathbb{P}(D\mid B_{3})\mathbb{P}(B_{3}).
\]

\end{example}

\section{Formula di Bayes}

Nelle stesse ipotesi della sezione precedente, per ogni $j=1,\ldots,n$ vale
\[
\mathbb{P}(B_{j}\mid A) = \frac{\mathbb{P}(A\mid B_{j})\mathbb{P}(B_{j})}
{\sum_{i=1}^{n}\mathbb{P}(A\mid B_{i})\mathbb{P}(B_{i})}.
\]


La formula di Bayes permette di aggiornare la probabilit\`{a} di uno scenario
$B_{j}$ dopo l'osservazione dell'evento $A$. In termini finanziari, essa
formalizza il passaggio da probabilit\`{a} a priori a probabilit\`{a} a posteriori.

Nel linguaggio del rischio di credito, $B_{j}$ pu\`{o} rappresentare uno stato
latente del debitore e $A$ un segnale osservabile, come un peggioramento dello
spread. La formula di Bayes consente allora di calcolare la probabilit\`{a}
aggiornata dello stato del debitore dopo il segnale.

\section{Esempio finanziario integrato}

Consideriamo un debitore che pu\`{o} trovarsi in uno di tre stati latenti:
\[
B_{1}=\text{solido},\qquad B_{2}=\text{fragile},\qquad B_{3}=\text{critico}.
\]
Supponiamo che tali stati formino una partizione di $\Omega$ e che le
probabilit\`{a} a priori siano
\[
\mathbb{P}(B_{1})=0.60,\qquad\mathbb{P}(B_{2})=0.30,\qquad\mathbb{P}%
(B_{3})=0.10.
\]
Si osserva un segnale negativo di mercato, indicato con $S$. Le
probabilit\`{a} del segnale condizionate allo stato del debitore sono
\[
\mathbb{P}(S\mid B_{1})=0.10,\qquad\mathbb{P}(S\mid B_{2})=0.40,\qquad
\mathbb{P}(S\mid B_{3})=0.80.
\]


La probabilit\`{a} non condizionata del segnale \`{e}
\[
\begin{aligned}
\Prob(S)
&=\Prob(S\mid B_1)\Prob(B_1)
+\Prob(S\mid B_2)\Prob(B_2)
+\Prob(S\mid B_3)\Prob(B_3)\\
&=0.10\cdot 0.60+0.40\cdot 0.30+0.80\cdot 0.10\\
&=0.26.
\end{aligned}
\]


La probabilit\`{a} a posteriori dello stato critico, dato il segnale negativo,
\`{e}
\[
\mathbb{P}(B_{3}\mid S)=\frac{\mathbb{P}(S\mid B_{3})\mathbb{P}(B_{3}%
)}{\mathbb{P}(S)}=\frac{0.80\cdot0.10}{0.26}\simeq0.3077.
\]
L'osservazione del segnale negativo aumenta quindi la probabilit\`{a} dello
stato critico dal $10\%$ a circa il $30.77\%$.

Supponiamo ora che la probabilit\`{a} di default condizionata allo stato sia
\[
\mathbb{P}(D\mid B_{1})=0.01,\qquad\mathbb{P}(D\mid B_{2})=0.05,\qquad
\mathbb{P}(D\mid B_{3})=0.20.
\]
Per calcolare la probabilit\`{a} di default dopo il segnale si determinano
prima le probabilit\`{a} posteriori degli stati:
\[
\mathbb{P}(B_{1}\mid S)=\frac{0.10\cdot0.60}{0.26}\simeq0.2308,
\]%
\[
\mathbb{P}(B_{2}\mid S)=\frac{0.40\cdot0.30}{0.26}\simeq0.4615,
\]%
\[
\mathbb{P}(B_{3}\mid S)=\frac{0.80\cdot0.10}{0.26}\simeq0.3077.
\]
La probabilit\`{a} di default condizionata al segnale \`{e} quindi
\[
\begin{aligned}
\Prob(D\mid S)
&=\Prob(D\mid B_1)\Prob(B_1\mid S)
+\Prob(D\mid B_2)\Prob(B_2\mid S)
+\Prob(D\mid B_3)\Prob(B_3\mid S)\\
&=0.01\cdot 0.2308+0.05\cdot 0.4615+0.20\cdot 0.3077\\
&\simeq 0.0869.
\end{aligned}
\]
Il segnale negativo porta quindi la probabilit\`{a} di default aggiornata a
circa l'$8.69\%$. L'esempio mostra come probabilit\`{a} totali e formula di
Bayes permettano di passare da un modello di scenari a una valutazione
aggiornata del rischio.

\section{Esercizi svolti}

\begin{exercise}
Sia
\[
\Omega=\{\omega_{1},\omega_{2},\omega_{3},\omega_{4}\}
\]
con probabilit\`{a} degli esiti elementari pari a $0.20,0.30,0.10,0.40$.
Siano
\[
A=\{\omega_{1},\omega_{2}\},\qquad B=\{\omega_{2},\omega_{4}\}.
\]
Calcolare $\mathbb{P}(A)$, $\mathbb{P}(B)$, $\mathbb{P}(A\cap B)$,
$\mathbb{P}(A\cup B)$ e $\mathbb{P}(A^{c})$.
\end{exercise}

\begin{solution}
Si ha
\[
\mathbb{P}(A)=0.20+0.30=0.50,
\]
\[
\mathbb{P}(B)=0.30+0.40=0.70,
\]
\[
A\cap B=\{\omega_{2}\}, \qquad\mathbb{P}(A\cap B)=0.30.
\]
Pertanto
\[
\mathbb{P}(A\cup B)=\mathbb{P}(A)+\mathbb{P}(B)-\mathbb{P}(A\cap B)
=0.50+0.70-0.30=0.90.
\]
Infine
\[
\mathbb{P}(A^{c})=1-\mathbb{P}(A)=0.50.
\]

\end{solution}

\begin{exercise}
Sia $D$ l'evento di default e sia $S$ un segnale negativo. Supponiamo che
\[
\mathbb{P}(D\cap S)=0.04,\qquad\mathbb{P}(S)=0.20.
\]
Calcolare $\mathbb{P}(D\mid S)$.
\end{exercise}

\begin{solution}
Per definizione di probabilit\`{a} condizionata,
\[
\mathbb{P}(D\mid S)=\frac{\mathbb{P}(D\cap S)}{\mathbb{P}(S)}=\frac
{0.04}{0.20}=0.20.
\]
La probabilit\`{a} di default condizionata al segnale negativo \`{e} quindi
pari al $20\%$.
\end{solution}

\section{Esercizi proposti}

\begin{exercise}
Si consideri uno spazio campionario con tre scenari di mercato: ribasso,
stabilit\`{a} e rialzo. Si definiscano due eventi finanziari rilevanti e si
calcolino unione, intersezione e complemento.
\end{exercise}

\begin{exercise}
Siano $A$ e $B$ due eventi con
\[
\mathbb{P}(A)=0.40,\qquad\mathbb{P}(B)=0.50,\qquad\mathbb{P}(A\cap B)=0.20.
\]
Verificare se $A$ e $B$ sono indipendenti.
\end{exercise}

\begin{exercise}
Un debitore pu\`{o} trovarsi in due stati, solido e fragile. Le
probabilit\`{a} a priori sono $0.70$ e $0.30$. La probabilit\`{a} di osservare
un segnale negativo \`{e} $0.10$ nello stato solido e $0.50$ nello stato
fragile. Calcolare la probabilit\`{a} a posteriori dello stato fragile dopo
l'osservazione del segnale.
\end{exercise}

\section{Sintesi finale}

Uno spazio di probabilit\`{a} $(\Omega,\mathcal{F},\mathbb{P})$ fornisce la
rappresentazione formale dell'incertezza. Lo spazio campionario $\Omega$
descrive gli esiti possibili; la sigma-algebra $\mathcal{F}$ descrive gli
eventi considerati dal modello; la misura $\mathbb{P}$ assegna probabilit\`{a}
agli eventi.

Le operazioni sugli eventi consentono di rappresentare eventi alternativi,
congiunti o complementari. La probabilit\`{a} condizionata formalizza
l'effetto dell'informazione sulla valutazione probabilistica. L'indipendenza
rappresenta il caso in cui tale informazione non modifica la probabilit\`{a}
dell'evento di interesse.

La formula delle probabilit\`{a} totali permette di decomporre una
probabilit\`{a} rispetto a una partizione di scenari. La formula di Bayes
permette di aggiornare le probabilit\`{a} degli scenari dopo l'osservazione di
un segnale. Questi strumenti costituiscono la base per la trattazione delle
variabili casuali, dei valori attesi condizionati e dei processi stocastici
nei capitoli successivi.

